\documentclass{article}
\usepackage{amsmath}
\usepackage{amssymb}
\usepackage{tikz}
\usepackage{booktabs}
\usepackage[margin=1in]{geometry}

\begin{document}

\section{The Model}

\subsection{Market Environment and Balance Sheets}

We consider a sequential economy operating in discrete time $t \in \{0, 1, 2, \dots\}$,
populated by a finite set of $N$ heterogeneous banks. Each bank $i$ operates under the
balance sheet identity
\begin{equation}
    L_{i} + C_{i} + R_{i} = D_{i} + E_{i}
\end{equation}
where $L_{i}$ are long-term assets (granted loans), $C_{i}$ is available lending capacity,
$R_{i}$ are reserves, $D_{i}$ are customer deposits, and $E_{i}$ is equity capital.
Time subscripts are omitted where no ambiguity arises.

%---------------------------------------------------------------
\subsection{Liquidity Shocks and Market Roles}

Each period, banks receive exogenous deposit shocks $\Delta D_{i,t}$. A positive realization
($\Delta D_{i,t} > 0$) generates excess liquidity, making bank $i$ a potential lender;
a negative realization ($\Delta D_{i,t} < 0$) creates a shortfall, making bank $i$ a
potential borrower. Lending capacity $C_{i,t}$ is thus determined by deposit dynamics.

%---------------------------------------------------------------
\subsection{Default and Bailout Probabilities}

The default probability of borrower $j$ is defined as a decreasing function of its
equity relative to the system maximum:
\begin{equation}
    p_{j} = 1 - \frac{E_{j}}{E_{\max}}
    \quad \Rightarrow \quad
    \begin{cases} p_{j} = 1 & \text{if } E_{j} = 0 \\ p_{j} = 0 & \text{if } E_{j} = E_{\max} \end{cases}
\end{equation}

The bailout probability of borrower $j$ scales with its relative asset size:
\begin{equation}
    b_{j} = \frac{A_{j,t-1}}{A_{\max,t-1}}
    \quad \Rightarrow \quad
    \begin{cases} b_{j} = 1 & \text{if } A_{j} = A_{\max} \\ b_{j} = 0 & \text{if } A_{j} = 0 \end{cases}
\end{equation}
$b_j$ is evaluated at $t-1$ to reflect that lenders observe the borrower's asset size
prior to the current period's shock, introducing a one-period information lag consistent
with the timing of interbank negotiations.\footnote{This lag also prevents simultaneity
between the shock that determines market roles and the size signal that governs bailout
expectations.}

%---------------------------------------------------------------
\subsection{Expected Loss and the Basel IRB Framework}

In the event of default, the lender's recovery depends on whether a state bailout occurs.
We map our parameters directly onto the Basel Internal Ratings-Based (IRB) framework,
where the expected loss (EL) of a credit exposure is a function of the Probability of
Default (PD), the Exposure at Default (EAD), and the Loss Given Default (LGD).

In our model, the PD is $p_j$ and the EAD is the loan size $L_{ij}$. The absolute
monetary loss conditional on default, denoted $\mathbb{E}(L \mid d)$ (which maps to
$\text{EAD} \times \text{LGD}$), is defined as:
\begin{equation}
    \mathbb{E}(L \mid d) = (1 - b_{j})(L_{ij} - \alpha A_{j}) + b_{j}(1 - \eta)L_{ij}
\end{equation}
Without a bailout, the lender recovers collateral $\alpha A_j$ and incurs a net loss
$L_{ij} - \alpha A_j$. Under a bailout, the state absorbs a fraction $\eta$ of the loan,
leaving a residual loss $(1-\eta)L_{ij}$. We impose $0 < \alpha \ll \eta < 1$, reflecting
that bailout recovery dominates collateral liquidation.

%---------------------------------------------------------------
\subsection{Internal Capital Constraint and Loan Supply}

Under the IRB framework, banks hold equity capital as a buffer against loss volatility.
Lender $i$ requires that the expected loss on any bilateral exposure not exceed a strict
fraction $\gamma$ of its equity base $E_i$, serving as an internal capital adequacy
constraint:\footnote{This formulation is consistent with the capital constraint
specifications of Nier et al.\ (2007) and Iori et al.\ (2006), and follows directly from
the balance sheet structure in Lenzu and Tedeschi (2012).}
\begin{equation}
    p_{j}\,\mathbb{E}(L \mid d) \leq \gamma E_{i}
\end{equation}
Setting the expected loss constraint to bind and solving for $L_{ij}$, subject to the
absolute lending capacity ceiling $C_i$, yields the optimal loan supply:
\begin{equation}
\boxed{
    L_{ij} = \min\!\left(
        \frac{\gamma E_{i} + p_{j}(1 - b_{j})\alpha A_{j}}{p_{j}(1 - b_{j}\eta)},\;
        C_{i}
    \right)
}
\end{equation}
By imposing a zero expected economic profit condition, we solve for the competitive
interest rate $r_{ij}$:
\begin{equation}
    \mathbb{E}(\pi) = (1 - p_{j})\,r_{ij}L_{ij} - p_{j}\,\mathbb{E}(L \mid d) - \kappa = 0
\end{equation}
\begin{equation}
\boxed{
    r_{ij} = \frac{p_{j}\,\mathbb{E}(L \mid d) + \kappa}{(1 - p_{j})\,L_{ij}}
}
\end{equation}
where $\kappa$ denotes exogenous per-loan operational costs. Note that $r_{ij} \to \infty$
as $p_j \to 1$: borrowers approaching absolute insolvency are priced out of the
market entirely.
%---------------------------------------------------------------
\subsection{Network Formation}

At the beginning of each period, borrowers revise their set of potential lending
agreements via a fitness mechanism. The fitness of lender $i$ is microfounded directly
from the loan supply equation. Since $L_{ij}$ is linear in $E_i$,
\begin{equation}
    \frac{\partial L_{ij}}{\partial E_{i}} = \frac{\gamma}{p_{j}(1 - b_{j}\eta)} > 0
\end{equation}
a borrower seeking to maximize credit access rationally targets the most capitalized
lender. Moreover, the TBTF amplification condition $\eta > p_j\alpha A_j /
(\gamma E_i + p_j \alpha A_j)$ is decreasing in $E_i$, so well-capitalized lenders
are precisely those for whom the bailout distortion is most operative. The fitness of
lender $i$ is therefore defined as its normalized equity ratio:
\begin{equation}
    \Phi_{i} = \frac{E_{i}}{E_{\max}}
\end{equation}
Borrower $j$ severs its link to lender $k$ and rewires to a randomly drawn lender $i$
with probability:
\begin{equation}
    P_{ij} = \frac{1}{1 + e^{-\beta(\Phi_{i} - \Phi_{k})}}
\end{equation}
where $\beta \geq 0$ governs the intensity of choice. Low values of $\beta$ yield
near-random topologies; high values amplify fitness differentials, concentrating
borrowing relationships on the most capitalized lenders.
%---------------------------------------------------------------





%------------------------------------------------------------------------------------
\subsection{Comparative Statics}

The interior loan supply $L_{ij}$ has the following properties.

\medskip
\noindent\textbf{(i) Effect of default risk.}
\begin{equation}
    \frac{\partial L_{ij}}{\partial p_{j}} = -\frac{\gamma E_{i}}{p_{j}^{2}(1 - b_{j}\eta)} < 0
\end{equation}
Loan supply is strictly decreasing in default probability, as higher risk tightens
the expected loss constraint.

\medskip
\noindent\textbf{(ii) Effect of bailout probability.}

Differentiating the interior solution with respect to $b_j$:
\begin{equation}
    \frac{\partial L_{ij}}{\partial b_{j}}
    = \frac{\eta\,\gamma E_{i} + p_{j}\alpha A_{j}(\eta - 1)}{p_{j}(1 - b_{j}\eta)^{2}}
\end{equation}
The denominator $p_j(1-b_j\eta)^2$ is strictly positive for $p_j > 0$. The sign of
the derivative is therefore determined entirely by the numerator. Setting it to be
positive:
\begin{align}
    \eta\,\gamma E_{i} + p_{j}\alpha A_{j}(\eta - 1) &> 0 \notag \\
    \eta\,\gamma E_{i} + p_{j}\alpha A_{j}\eta - p_{j}\alpha A_{j} &> 0 \notag \\
    \eta(\gamma E_{i} + p_{j}\alpha A_{j}) &> p_{j}\alpha A_{j}
\end{align}
which yields the condition:
\begin{equation}
\boxed{
    \frac{\partial L_{ij}}{\partial b_{j}} > 0
    \iff
    \eta > \frac{p_{j}\alpha A_{j}}{\gamma E_{i} + p_{j}\alpha A_{j}}
}
\end{equation}
The right-hand side is bounded in $(0,1)$ whenever $\gamma E_i > 0$. Under the
maintained assumption $\alpha \ll \eta$, the right-hand side is small and the condition
is satisfied generically. The intuition is direct: a higher bailout probability relaxes
the loss-given-default, which loosens the expected loss constraint and allows the lender
to extend more credit. The strength of this effect increases with $\eta$ and diminishes
with $\alpha$ and $p_j$.

\medskip
\noindent\textit{Boundary case $\gamma E_i = 0$.} If the lender holds no equity buffer,
the expected loss constraint collapses to $p_j \mathbb{E}(L \mid d) = 0$, so the
loan is backed purely by collateral:
\begin{equation}
    L_{ij} = \frac{(1 - b_{j})\,\alpha A_{j}}{1 - b_{j}\eta} \leq \alpha A_{j}
\end{equation}
In this case the numerator of $\partial L_{ij}/\partial b_j$ becomes
$p_j \alpha A_j (\eta - 1) < 0$, so the derivative reverses sign. The boundary
cases are summarised in Table~\ref{tab:boundary}.

\begin{table}[h]
\centering
\caption{Boundary cases for loan supply and interest rate}
\label{tab:boundary}
\begin{tabular}{lll}
\toprule
Condition & $L_{ij}$ & $r_{ij}$ \\
\midrule
$b_j = 0$ & $\dfrac{\gamma E_i + p_j \alpha A_j}{p_j}$ & $\dfrac{p_j(L_{ij} - \alpha A_j) + \kappa}{(1-p_j)L_{ij}}$ \\[10pt]
$b_j = 1$ & $\dfrac{\gamma E_i}{p_j(1-\eta)}$ & $\dfrac{p_j(1-\eta)L_{ij} + \kappa}{(1-p_j)L_{ij}}$ \\[10pt]
$p_j = 0$ & $C_i$ & $\dfrac{\kappa}{C_i}$ \\[10pt]
$p_j = 1$ & $\dfrac{\gamma E_i + (1-b_j)\alpha A_j}{1 - b_j\eta}$ & $\to \infty$ \\
\bottomrule
\end{tabular}
\end{table}

\end{document}